\subsection{Conditional exclusion from a stable generator and a scalar deformation law}
\label{sec:pauli-stability-selection}

One can select exclusion without assuming it, provided the quantization
class and the relation between occupation and dynamics are specified.
The following theorem permits a real scalar deformation parameter rather
than declaring either canonical commutation or anticommutation relations.
Its phase-generator and stability hypotheses are additional physical
inputs; a finite spin representation alone supplies neither.

\begin{theorem}[Stability selects exclusion within a scalar deformation class]
\label{thm:pauli-stability-selection}
Let $\mathcal K$ be a finite-dimensional complex one-particle Hilbert
space. On a positive Hilbert space $\mathcal F$, let
$v\mapsto a^\dagger(v)$ be complex-linear, with
$a(v)$ its adjoint on a common dense invariant domain $\mathcal D$.
Assume $\mathcal D$ is invariant under these operators and a self-adjoint
Hamiltonian $H$, and is contained in their domains. Supply:
\begin{enumerate}
 \item a common real number $q$ such that, for every unit $v$,
 \[
  a(v)a^\dagger(v)-q\,a^\dagger(v)a(v)=I
 \]
 on $\mathcal D$;
 \item the exact occupation/phase law, for every unit $v$,
 \[
  N_v=a^\dagger(v)a(v),\qquad [N_v,a^\dagger(v)]=a^\dagger(v);
 \]
 \item a normalized $\Omega\in\mathcal D$ with
 $a(v)\Omega=0$ for all $v$ and $H\Omega=E_0\Omega$;
 \item a unit mode $v_-$ and $\kappa>0$ such that
 \[
  [H,a^\dagger(v_-)]=-\kappa a^\dagger(v_-)
 \]
 on $\mathcal D$, and a finite uniform lower bound for $H$.
\end{enumerate}
Then $q=-1$, and on $\mathcal D$,
\[
 \{a(v),a^\dagger(w)\}=\langle v,w\rangle I,
 \qquad \{a^\dagger(v),a^\dagger(w)\}=0,
 \qquad \{a(v),a(w)\}=0.
\]
Here the inner product is conjugate-linear in its first argument. In
particular $(a^\dagger(v))^2=0$ for every mode $v$. The operators extend
boundedly to $\mathcal F$, where the same relations hold.
The reference vacuum $\Omega$ is not assumed to be a ground state.
\end{theorem}
\begin{proof}
Fix a unit mode, abbreviate its operators to $a,a^\dagger,N$, and work
throughout on the common invariant domain. The deformation law gives
\[
 Na^\dagger=a^\dagger+q\,a^\dagger N.
\]
Comparison with the occupation/phase law yields
$(q-1)a^\dagger N=0$. The vector $|1\rangle=a^\dagger\Omega$
has norm one and $N|1\rangle=|1\rangle$. If $q\ne1$, it follows that
$a^\dagger|1\rangle=0$. Taking its squared norm using the deformation
law gives $0=1+q$. Thus $q$ is either $1$ or $-1$.

If $q=1$, induction using $aa^\dagger=I+a^\dagger a$ gives
\[
 |n\rangle=(a^\dagger)^n\Omega,\qquad
 \|\,|n\rangle\|^2=n!,\qquad N|n\rangle=n|n\rangle
\]
for every integer $n\geq0$. Apply this to $v_-$. The Hamiltonian
commutator and the vacuum eigenvalue imply
$H|n\rangle=(E_0-n\kappa)|n\rangle$. All these vectors are nonzero
and belong to $\mathcal D$, contradicting the lower bound. This is an
arbitrary-$n$ argument, not a finite occupation scan. Hence $q=-1$.

For any unit mode, the phase identity then gives $a^\dagger N=0$,
whereas the deformation law gives $Na^\dagger=a^\dagger$. Consequently
\[
 0=a^\dagger Na^\dagger=(a^\dagger)^2.
\]
Linearity and scaling extend this square-zero identity to every $v$,
including $v=0$. Expanding it at $v+w$ proves creation
anticommutation, and taking adjoints on the dense invariant domain proves
annihilation anticommutation. The diagonal relation is
$\{a(v),a^\dagger(v)\}=\|v\|^2I$ for every $v$. Polarization at
$v+w$ and $v+iw$ gives the mixed relation. Finally, for
$\psi\in\mathcal D$,
\[
 \|a(v)\psi\|^2+\|a^\dagger(v)\psi\|^2
   =\|v\|^2\|\psi\|^2.
\]
Both operators therefore have bounded extensions, and density extends
their identities to all of $\mathcal F$.
\end{proof}

\begin{corollary}[Finite-dimensional alternative to the stability hypothesis]
\label{cor:pauli-finite-capacity-selection}
Retain the theorem's common scalar deformation law, linear mode fields,
exact occupation/phase law and normalized annihilation vacuum, with
$\mathcal K\ne\{0\}$. If $\mathcal F$ itself is nonzero and
finite-dimensional, the same exclusion
and CAR conclusions hold without assuming a Hamiltonian or a negative
ladder.
\end{corollary}
\begin{proof}
A dense linear domain in finite dimension is the whole space. The
vacuum argument gives $q\in\{1,-1\}$. If $q=1$, the trace of
$aa^\dagger-a^\dagger a=I$ is both zero and
$\dim\mathcal F>0$, a contradiction. The square-zero and polarization
arguments for $q=-1$ require no Hamiltonian.
\end{proof}

This alternative requires finite dimension of the Hilbert space carrying
the full creation and annihilation algebra. A finite one-particle mode
space alone is insufficient.
A finite observer response space does not by itself identify that
operator representation with physical matter. Applying the corollary to
observer finite capacity requires precisely this additional identification,
as well as the stated deformation and phase laws.

The negative-mode hypothesis has a concrete one-particle instance in the
supplied finite matter hopping action. For a unitary spin-link matrix $K$
and $\kappa>0$, its edge matrix is
\[
 h_e=\kappa\begin{pmatrix}0&K\\K^\dagger&0\end{pmatrix},
 \qquad h_e^2=\kappa^2I.
\]
For any unit spin vector $\chi$, the vector
$v_-=(-K\chi,\chi)/\sqrt2$ is normalized and satisfies
$h_ev_-=-\kappa v_-$. The quantum dynamical law
$[H,a^\dagger(v)]=a^\dagger(h_ev)$ would therefore supply the
required ladder. This law is a hypothesis linking the one-particle
action to its quantization. A one-particle matrix, a measured unitary
step, or the existence of a spin lift does not establish that law.

There is a finite nonvacuity witness. In ordinary fermionic Fock space
over the edge modes, take the empty vacuum and
$H=d\Gamma(h_e)$ with its normal ordering fixed at that vacuum.
The negative-energy modes can each be occupied once, and the ground
energy is the sum of all negative one-particle eigenvalues. Filling those
modes gives the ground state; particles in positive modes and holes in
filled negative modes both have excitation energy $\kappa$.
For the two-spin edge with fifteen internal channels of the supplied
matter census, there are thirty negative modes. Its ground energy is
$-30\kappa$. With
\[
 (d_s)=(6,3,3,2,1),\qquad (q_s)=(1,-4,2,-3,6),
 \qquad \sum_s d_sq_s=0,
\]
the filled ground state has total hypercharge zero. Subtracting its
constant energy gives a positive excitation Hamiltonian and preserves
the field commutator. The reference empty vacuum and the filled ground
state are different vectors.

For comparison, canonical bosonic Fock space on the same fixed edge
action admits occupation $n$ in one negative mode of every internal
channel. These states have total hypercharge zero and energy
$-15n\kappa$. Thus even the globally neutral sector has no uniform
lower bound. Adding a chemical potential proportional to total
hypercharge cannot repair this sequence. This neutrality statement does
not impose local quantum Gauss constraints or nonabelian singlet
constraints, and is not a claim about a gauge-invariant physical Hilbert
space.

The hypotheses exclude several concrete alternatives. At fixed total
particle number the bosonic Hamiltonian is bounded below, so stability
across all occupation sectors is essential. Replacing the one-particle
law by $h_e+cI$ with $c\geq\kappa$ stabilizes its full bosonic Fock
Hamiltonian: its second-quantized change is $cN_{\rm total}$, not a
constant vacuum-energy subtraction. Positive number interactions can
also bound the energy, while changing the negative ladder. Neither
change is ruled out by the theorem. A positive spectral replacement
such as $|h_e|$ likewise changes the supplied dynamics.

Even the phase law and stability together do not force exclusion.
On a three-dimensional Hilbert space with basis $|0\rangle,|1\rangle,
|2\rangle$, set
\[
 a^\dagger|0\rangle=|1\rangle,\qquad
 a^\dagger|1\rangle=\sqrt2|2\rangle,\qquad
 a^\dagger|2\rangle=0.
\]
Then $N=a^\dagger a=\operatorname{diag}(0,1,2)$ satisfies
$[N,a^\dagger]=a^\dagger$, and $H=-\kappa N$ is bounded below,
but $(a^\dagger)^2|0\rangle\ne0$. This model violates the scalar
deformation hypothesis: $aa^\dagger=\operatorname{diag}(1,2,0)$,
whose middle entry requires $q=1$ while the top entry then contradicts
$aa^\dagger-qN=I$. A hard occupation cutoff or a more general exchange
algebra is therefore a genuine alternative outside the stated class.

The theorem is a conditional Pauli-exclusion derivation within a specified
quantization class. It is insensitive to spin, so it cannot attach
half-integer spin to fermionic statistics or exclude a different
quantization class for a different field. The common scalar deformation
law, linear field structure, exact occupation/phase law, unrestricted
ladder domain and stable physical generator remain supplied. The finite
filled-sea construction supplies no continuum vacuum or relativistic
locality theorem. Deriving these inputs and selecting the physical
matter action from observer consistency remain separate requirements.
